\section{Introduction}
\label{sec:introduction}

\subsection{Urban heat as a decision problem}

Urban populations are exposed to temperatures systematically higher than those
of their rural surroundings, and the excess is not uniform across a city. The
distinction between the surface, canopy-layer, and boundary-layer urban heat
island, established by \textcite{oke1976}, remains the necessary starting point
for any serious treatment of the subject, because the three are different
physical quantities measured in different ways and they do not agree in
magnitude. Within a single city, daytime near-surface air temperature can vary
by several degrees over short distances: \textcite{ziter2019} measured a mean
daytime intra-urban air temperature range of \SI{3.5}{\degreeCelsius} (spanning
\SIrange{1.1}{5.7}{\degreeCelsius}) and a mean nighttime range of
\SI{2.1}{\degreeCelsius} (\SIrange{1.2}{3.0}{\degreeCelsius}) across a single
urban area. Variation of this magnitude, over distances a person walks in
minutes, is what makes urban heat a spatial planning problem rather than only a
meteorological one.

The consequences of that variation are distributed unevenly across the
population. \textcite{reid2009} reduced ten candidate vulnerability variables ---
six demographic, two describing household air-conditioning, one describing
vegetation cover, and one describing diabetes prevalence --- to four factors
explaining more than \SI{75}{\percent} of total variance: social and
environmental vulnerability, social isolation, air-conditioning prevalence, and
the proportion of elderly and diabetic residents. Two people experiencing the
same air temperature do not face the same risk. Any instrument that ranks sites
for intervention on physical heat alone is answering a narrower question than
the one cities are actually asking.

\subsection{Nature-based solutions as a cooling response}

Vegetation and water cool urban environments through shade, evapotranspiration,
and evaporation, and the empirical evidence for the effect is well established
at the scale of individual green sites. \textcite{bowler2010}, restricting their
meta-analysis to empirical studies, found that an urban park was on average
\SI{0.94}{\degreeCelsius} cooler during the day than its non-green
surroundings, with larger parks and parks containing trees cooler than smaller
or open ones. \textcite{ziter2019} showed that the relationship between canopy
cover and daytime air temperature is nonlinear, with the greatest cooling
occurring above approximately \SI{40}{\percent} canopy cover, and that the
effect is scale-dependent, peaking at radii of \SIrange{60}{90}{\metre} ---
roughly the scale of a city block. The same study found that impervious cover
raises temperature approximately linearly and that its warming effect persists
overnight, where canopy cooling largely does not.

Aggregating across measures and climates, \textcite{keravec2026} conducted an
umbrella review of 64 systematic reviews with a multilevel meta-regression of
pedestrian-level daytime air temperature reductions. That synthesis is the
single most useful cross-typology anchor currently available, and it carries a
finding with direct consequences for tool design: the type of intervention
alone explains only \SI{12}{\percent} of the variance in observed cooling
effectiveness, rising to \SI{30}{\percent} when climate zone is added and to
\SI{78}{\percent} once the K\"oppen--Geiger climate subzone is included. An
instrument that reports one global cooling figure per intervention type is
therefore discarding the majority of the explanatory signal in the evidence
base.

\subsection{The decision gap}

A city preparing a heat adaptation programme must choose among many more
candidate sites and interventions than it can study in detail. In practice it
has two options, and neither serves the triage step.

The first is unstructured judgment, often mediated by a spreadsheet. It is fast
and cheap, and it draws on genuine local expertise. Its weakness is not that it
is wrong but that it is not \emph{comparable}: two officers assessing two sites
apply different implicit weights, cite no evidence that a third party can
check, and produce numbers that cannot be defended to a funder or a public
inquiry. Nothing about the process records why one site outranked another.

The second is detailed simulation --- microclimate models of the ENVI-met
class, computational fluid dynamics, or building energy simulation. These are
the appropriate instruments for design, and they are rigorous. The work of
\textcite{zolch2016} illustrates both their value and their cost: a micro-scale
ENVI-met evaluation of heat mitigation measures which established, among other
things, that trees achieved a mean physiologically equivalent temperature (PET)
reduction of \SI{13}{\percent} at pedestrian level, that green fa\c{c}ades
achieved \SIrange{5}{10}{\percent}, that green roof effects at pedestrian level
were negligible, and --- a finding of general importance --- that increasing
green cover did not correspond linearly to PET reduction, because placement
matters more than quantity. Producing such results requires geometry, material
properties, meteorological forcing, and specialist time. A city cannot run them
across two hundred candidate street segments before deciding which twenty to
investigate.

The gap between the two is the screening step: a structured, repeatable,
evidence-referenced first pass that narrows a long list to a short one and
records why. That is the step this methodology addresses.

\subsection{Contribution and structure of this paper}

This paper specifies, in full, the methodology of \tool{} at methodology
version \methver{}. It is written to be read independently of the software, by
reviewers who wish to evaluate the scientific basis of the tool's outputs
before deciding whether to trust, adopt, or contest them. Its contributions
are four.

\begin{enumerate}
  \item \textbf{A complete, auditable specification.} Every formula the engine
  applies is stated and typeset here, every weight is declared with its
  justification or with an explicit statement that it is a value choice, and
  every performance value is traced to a cited source whose verification
  status is disclosed (\Cref{app:verification}).

  \item \textbf{An evidence-derived library with its evidence classes made
  explicit.} \num{110} curated catalogue entries each inherit one of \num{18}
  \emph{cooling archetypes}, and the archetype --- not the entry --- carries
  the envelope, the base score, the evidence rating, and the citations. The
  reason is stated plainly rather than finessed: solution-specific cooling
  literature does not exist for \num{110} intervention types, so an entry
  inherits a \emph{cited} envelope and every result names the evidence class it
  inherited from (\Cref{sec:archetypes}). Calibration decisions that materially
  lower the tool's headline numbers --- a downward correction for modelling
  bias, the refusal of super-additivity for multi-intervention packages, the
  deliberate non-adoption of newer and larger published figures, and
  conservative treatment where sources conflict --- are presented for scrutiny
  in \Cref{sec:calibration}.

  \item \textbf{A discipline of honest uncertainty.} All quantitative outputs
  are ranges; every output block carries its own confidence rating; missing
  data reduces confidence rather than being silently imputed; and the
  temperature-envelope clipping rule prevents favourable site conditions from
  generating claims that exceed published evidence (\Cref{sec:cooling}).

  \item \textbf{Reproducibility as an engineering property.} The methodology
  lives in versioned configuration files rather than in code, and its evidence
  rules are enforced by the build: an uncited performance value or a citation
  pointing at an unrecorded source fails continuous integration
  (\Cref{sec:implementation}). This is what allows the claims in this paper to
  be checked rather than merely believed.
\end{enumerate}

\Cref{sec:scope} states what the tool is and is not, and the design commitments
that constrain it. \Cref{sec:related} positions it against published work and
argues for the standards by which a screening instrument should be judged.
\Cref{sec:framework} presents the three-layer conceptual structure.
\Cref{sec:inputs,sec:typologies,sec:formulas} specify the inputs, the typology
library, and every scoring formula.
\Cref{sec:uncertainty,sec:sensitivity} cover uncertainty treatment and the
executed sensitivity analysis. \Cref{sec:implementation} describes the
implementation and the machine-enforced evidence rules.
\Cref{sec:worked-example} works one illustrative assessment through by hand.
\Cref{sec:limitations,sec:governance} state the limitations and the governance
process. Appendices reproduce the archetype library and the catalogue entries
that inherit from it, the complete weight tables, the input field reference, a
glossary, and the source verification register.
